\chapter*{Remerciements}
\addcontentsline{toc}{chapter}{Remerciements}

Ce travail a été réalisé dans le cadre d'un stage effectué au sein du Comité National de l'Initiative pour la Transparence dans les Industries Extractives du Sénégal (ITIE Sénégal). Nous tenons à remercier l'ensemble de l'équipe technique du Comité pour son accueil, la mise à disposition de la documentation sur le secteur extractif sénégalais et les échanges qui ont nourri la réflexion présentée dans ce rapport.

Nos remerciements s'adressent également au corps enseignant de l'École Nationale de la Statistique et de l'Analyse Économique de Dakar (ENSAE Dakar), dont la formation en modélisation économique et finances publiques a constitué le socle méthodologique de ce travail.

Enfin, ce travail s'appuie largement sur les outils et publications du \emph{Natural Resource Governance Institute} (NRGI) et du Fonds Monétaire International (FMI), ainsi que sur les travaux académiques portant sur l'application du modèle FARI au secteur minier ouest-africain, en particulier le mémoire de \citet{kourouma2019} consacré au cas guinéen, qui a servi de point de repère méthodologique pour la présente étude appliquée au cas sénégalais.
